% !Mode:: "TeX:UTF-8"
% !Mode:: "TeX:UTF-8"
\newcommand{\thepage}{} % to silence a hyperref warning about pdfpagelabels
\RequirePackage[]{silence}
\WarningsOff[hyperref]
\PassOptionsToPackage{log-declarations=false}{xparse}
%\RequirePackage[l2tabu, orthodox]{nag}
\RequirePackage{luatex85} % to make \pdf...= commands work with LuaTeX
\documentclass[aspectratio=169,t]{beamer}
\usepackage{ifthen}
\usepackage[T1]{fontenc}
\usepackage{iftex}
\ifLuaTeX
 \usepackage{fontspec}
 \setmainfont[Ligatures=TeX]{Source Sans Pro}
\else
 \usepackage[utf8]{inputenc}
\fi
\usepackage{german}
\usepackage{newunicodechar}
\newunicodechar{€}{\texteuro}
\newunicodechar{ß}{\ss}
\renewcommand{\SS}{\iffontchar\font"1E9E \symbol{"1E9E}\else SS\fi}
\usepackage{ellipsis}
\usepackage{microtype}
\usepackage{listings}
\usepackage{sourcecodepro}
\lstset{language=Python, basicstyle=\ttfamily}
\usepackage{amsfonts}
\usepackage{amssymb}
\usepackage{amsthm}
\usepackage{fancyhdr}
\usepackage{graphicx}

\newtheorem{proposition}[theorem]{Proposition}

\DeclareMathOperator*{\argmax}{argmax}
\DeclareMathOperator{\supp}{supp}
\newcommand{\E}{\ensuremath{\mathbb{E}}}
\newcommand{\V}{\ensuremath{\mathbb{V}}}

\setbeamercolor{emph}{fg=purple}
\renewcommand<>{\emph}[1]{%
  {\usebeamercolor[fg]{emph}\only#2{\itshape}#1}%
}



\hypersetup{draft=false,colorlinks=false,pdfpagemode=UseOutlines,plainpages=false,
            breaklinks=true,bookmarksnumbered=true,pdfborder=0 0 0,
            pdftitle={Data Science SoSe 2026 HHU},
            pdfauthor={Konrad Voelkel},
            pdfsubject={data science}}

%\documentclass[aspectratio=169,t]{beamer}
\usetheme{hhubeamer} % put the hhu-beamer-theme in /usr/share/texmf/tex/latex/beamer/themes/

% -------------------- Code listings config --------------------
% Define custom colors for syntax highlighting
\definecolor{codegreen}{rgb}{0,0.6,0}
\definecolor{codegray}{rgb}{0.5,0.5,0.5}
\definecolor{codepurple}{rgb}{0.58,0,0.82}
\definecolor{backcolour}{rgb}{0.96,0.96,0.94} % Very light gray/beige
% Configure the listings package
\lstset{
    language=Python,
    backgroundcolor=\color{backcolour},   
    commentstyle=\color{codegreen}\itshape,
    keywordstyle=\color{blue}\bfseries,
    numberstyle=\tiny\color{codegray},
    stringstyle=\color{codepurple},
    basicstyle=\ttfamily\footnotesize, % Use a monospaced font
    breakatwhitespace=false,         
    breaklines=true,      
             % Enable automatic line breaking
    captionpos=b,                      % Caption at the bottom
    keepspaces=true,                 
    numbers=left,                      % Line numbers on the left
    numbersep=10pt,                    % Distance of line numbers from the code
    showspaces=false,                
    showstringspaces=false,            % Don't show weird underscores in strings
    showtabs=false,                  
    tabsize=4,
    frame=single,                      % Add a subtle border around the code
    rulecolor=\color{black!20},        % Border color
  }

\event[]{}
\author[Sommersemester 2026]{Dr. Konrad Völkel, Sommersemester 2026}

%\title{Data Science, VL 1}
%\date{16.04.2026}


%%% Local Variables:
%%% mode: latex
%%% TeX-master: t
%%% End:


\title{Data Science, VL 20}
\subtitle{Lineare Regression}
\date{30.06.2026}

\begin{document}

\begin{frame}
  \titlepage
\end{frame}

\begin{frame}[separator]
  \frametitle{Lineare Regression}
\end{frame}

\begin{frame}
\frametitle{Beispiel zur Linearen Regression}
Wir haben bereits ein Beispiel gesehen: Blütenblattbreite aus Blütenblattlänge.

\vspace{0.5cm}

\uncover<2->{
  Weiteres Beispiel: Mit einem Datensatz von Häusern in einer australischen Großstadt wollen wir den Hauspreis $Y$ vorhersagen aus der Entfernung zum Stadtzentrum $X_1$ und dem Alter des Hauses $X_2$.
  }

\uncover<3->{
Das lineare Modell ist $Y = \beta_1 X_1 + \beta_2 X_2 + \beta_0 + \epsilon$.
}

\uncover<4->{
  Achtung: In der Informatik / im maschinellen Lernen üblich: \\
  $f(x_1,x_2;\beta_0,\beta_1,\beta_2) := \beta_1 x_2 + \beta_2 x_2 + \beta_0$ wird \alert{linear} genannt,
  obwohl nur $(x_1,x_2) \mapsto \beta_1 x_1 + \beta_2 x_2$ wirklich linear ist; \\
  $f$ heißt in der Mathematik \alert{affin-lineare Funktion}.
}

\vspace{0.5cm}
\uncover<5->{
  Wie finden wir nun 'gute' Parameter $\beta$ und was ist 'gut' für uns?
}

\end{frame}

\begin{frame}
\frametitle{Lineare Regression, formal}
% hier auch beide Varianten, mit und ohne Bias.
Ein Regressionsmodell der Form $Y = f(X;w) + \epsilon$ mit $f(x;w) = w^T x$ heißt\\ \alert{lineare Regression}.
\uncover<2->{Die Parameter $w$ heißen auch \alert{Gewichte} (weights).}

\vspace{0.5cm}

\uncover<3->{
Ein Regressionsmodell der Form $Y = f(X;w, w_0) + \epsilon$ mit $f(x;w) = w^T x + w_0$ heißt\\ \alert{lineare Regression} mit \alert{Bias} $w_0$.
}

\vspace{1cm}

\uncover<4->{
  Wenn wir statt $(X_1,\dots,X_k)$ als Eingabe $X_+ := (1,X_1,\dots,X_k)$ verwenden und nicht $w = (w_1,\dots,w_k)$ sondern $w_+ := (w_0,w_1,\dots,w_k)$
  dann ist $(w_+)^T X_+ = w^T X + w_0$.

  \vspace{0.5cm}

  Daher ist jede lineare Regression mit Bias auch eine sehr spezielle lineare Regression ohne Bias.
  }

\end{frame}

\begin{frame}[fragile]
  \frametitle{Lineare Regression mit Scipy.optimize}
  
\begin{lstlisting}[language=Python]
from scipy.optimize import curve_fit

def f(x, beta_0, beta_1):
    return beta_1*x + beta_0

params, covariance_matrix = curve_fit(f, x, y)
\end{lstlisting}

\uncover<2->{
  Die Variable \texttt{params} ist das 2-Tupel $\beta$.

  \vspace{0.5cm}

  Aus der Dokumentation von \texttt{curve\_fit}:
  
  Use non-linear least squares to fit a function, f, to data.

Assumes \texttt{ydata = f(xdata, *params) + eps}.
}
\end{frame}


\begin{frame}
\frametitle{Lineare Regression als Schätzer}
\begin{definition}
  Wenn man in einem Regressionsmodell die \alert{Residuenquadratsumme}
  $\sum_{i=1}^N {\left(y_i - f(x_i;\beta)\right)}^2$ minimiert,
  spricht man von einem Schätzer für $\beta$, der als \alert{kleinste-Quadrate-Schätzer} (least squares) bekannt ist.
\end{definition}

\uncover<2->{
Der kleinste-Quadrate-Schätzer ist ein Maximum-Likelihood-Schätzer.
}
\end{frame}


\begin{frame}
\frametitle{Kleinste Quadrate}
Wenn wir zu Punkten $(x_i,y_i)$ in der Ebene eine Gerade berechnen,
die den gemeinsamen Abstand zu allen Punkten minimiert (eine sogenannte Ausgleichsgerade),
so ist dies genau eine Regression mit $f(x, \beta) = \beta_0 + \beta_1 x$.
Die Regressionsgleichung $Y_i = f(X_i, \beta)$ wird zur Geradengleichung
$y_i = \beta_0 + \beta_1 x_i$
wobei $\beta_0$ der y-Achsenabschnitt (bei $x_i=0$) und $\beta_1$ die Steigung der Geraden ist.

\vspace{0.5cm}

\uncover<2->{
Die kleinste-Quadrate-Methode ist nun, $\beta_0$ und $\beta_1$ so zu schätzen,
dass $\sum_{i=1}^N {\left(y_i - \beta_1x_i - \beta_0  \right)}^2$ minimiert wird.
}
\end{frame}

\begin{frame}
\frametitle{Kleinste Quadrate von Hand, pt. 1}
Wenn $\beta_0 = 0$ ist (die Ausgleichsgerade durch den Ursprung verläuft), ist $\beta_1$ leicht zu schätzen:
Wir leiten die Residuenquadratsumme nach $\beta_1$ ab und finden Nullstellen.

\uncover<2->{
\begin{eqnarray*}
0 &= \sum_{i=1}^N 2{\left(y_i - \beta_1 x_i  \right)}{\left(- x_i  \right)} \\
  &= \sum_{i=1}^N 2 \left( \beta_1 x_i^2 - y_i x_i \right) \\
  &= 2\left(\beta_1 \sum_{i=1}^N x_i^2 - \sum_{i=1}^N x_iy_i \right)
\end{eqnarray*}

also ist $\beta_1 = \dfrac{\sum_{i=1}^N x_iy_i}{\sum_{i=1}^N x_i^2}$.
}
\end{frame}

\begin{frame}
\frametitle{Kleinste Quadrate von Hand, pt. 2}
Im Allgemeinen lässt sich $\beta_0$ berechnen mit der Formel

$$
\beta_0 = \texttt{mean}(y) - \beta_1\texttt{mean}(x).
$$

\vspace{0.5cm}

\uncover<2->{
Um also zunächst $\beta_1$ im Allgemeinen zu bestimmen, verschiebt man die Daten um die Mittelwerte, also ersetzt $x_i$ durch $x_i' = x_i - \texttt{mean}(x)$ und $y_i$ durch $y_i' = y_i - \texttt{mean}(y)$. Damit erhält man die Formel

$$
\beta_1 = \dfrac{\sum_{i=1}^N (x_i - \texttt{mean}(x))(y_i - \texttt{mean}(y))}{\sum_{i=1}^N (x_i - \texttt{mean}(x))^2}
$$
}
\end{frame}

\begin{frame}[fragile]
  \frametitle{Lineare Regression mit Scipy.Stats}
  Statt curve fitting lässt sich auch direkt lineare Regression verwenden:
  
\begin{lstlisting}
from scipy.stats import linregress

linregress_result = linregress(x, y)
slope, intercept, r, p, se = linregress_result
\end{lstlisting}

\uncover<2->{
  Dabei ist \texttt{slope} die Steigung, also $\beta_1$ und \texttt{intercept} der y-intercept, der y-Achsenabschnitt der Geraden, also der Bias $\beta_0$.
  Weiterhin sind \texttt{r} der Korrelationskoeffizient, \texttt{p} ein $p$-Wert und \texttt{se} der quadrierte Fehler (squared error), also das mittlere Residuenquadrat.
}
\end{frame}


\begin{frame}
\frametitle{Achtung: Verwechslungsgefahr}
% Gerade maximaler Varianz != Regressionsgerade

Die Regressionsgerade (=Ausgleichsgerade) ist \alert{nicht} die Richtung maximaler Varianz!

\uncover<2->{
  Zeichnet man die Abstände ein, die in beiden Fällen gemessen werden, sieht man den Unterschied:

  Bei der Regressionsgerade minimiert man die (Quadrate der) Residuen, die parallel zur y-Achse verlaufen,
  bei der Richtung maximaler Varianz minimiert man die Abstände zur Geraden, die senkrecht zur Geraden verlaufen.

  (Wir betrachten später noch mehr Bilder im Notebook).
}
\end{frame}

\begin{frame}
\frametitle{Multilineare Regression}

Wenn wir nur einen univariaten Regressor $X = (X_1)$ haben,
also Gewicht $w = w_1$ und Bias $w_0$, sprechen wir von einer \alert{einfachen linearen Regression}.

\uncover<2->{
Wenn der Regressor multivariat ist, also $X = (X_1,\dots,X_k)$ und damit das Gewicht ein Gewichtsvektor $w = (w_1,\dots,w_k)$ (mit Bias $w_0$) ist, sprechen wir von einer \alert{multilinearen Regression}.
}

\vspace{1cm}

\uncover<3->{
  Während   \texttt{scipy.stats.linregress} nur einfache lineare Regression unterstützt,
  bietet Scikit-Learn auch multilineare Regression an mit der Klasse \texttt{sklearn.linear\_model.LinearRegression}.
  }
\end{frame}


\begin{frame}[separator]
  \frametitle{Lineare Regression mit Basisfunktionserweiterung}
\end{frame}

\begin{frame}
\frametitle{Beispiel zur Basisfunktionserweiterung}
Gegeben bivariate Daten $(X_1,X_2)$ wie z.B. Blütenblattlänge und -breite,
könnten wir auch einen Zusammenhang zwischen Kelchblattlänge (sepal length) $Y$ und der Blütenblattfläche (petal leaf area) vermuten.

\uncover<2->{
Wir bilden ein neues Feature $X_3 := X_1 X_2$ (das Produkt) und betrachten die lineare Regression $Y = w_1 X_3 + w_0 + \epsilon$.
}
\end{frame}

\begin{frame}
\frametitle{Basisfunktionserweiterung}
Wir haben die Funktion $\phi(X_1,X_2) := X_1 X_2$ auf die ursprünglichen Daten angewendet.

Das ist eine \alert{Basisfunktionserweiterung} mit $\phi$.

\uncover<2->{
  \begin{definition}
    Die lineare Regression mit Regressor $\phi(X_1,\dots,X_k)$ wird lineare Regression mit Regressor $X_1,\dots,X_k$ und Basisfunktionserweiterung mit $\phi$ genannt.
  \end{definition}
}
\end{frame}

\begin{frame}
  \frametitle{Noch mehr Beispiele für Basisfunktionserweiterungen}

  Wir kennen bereits den Trick $\phi(X) = (1,X)$ bzw. $\phi(X_1,\dots,X_k) = (1,X_1,\dots,X_k)$.
  
  Eine lineare Regression mit dieser Basisfunktionserweiterung ist wie eine lineare Regression mit Bias-Term $w_0$.
\end{frame}

\begin{frame}
\frametitle{Polynomiale Regression}
Erweiterung mit $\phi(X) = (1,X,X^2)$ nennt man \alert{quadratische Regression},
allgemeiner $\phi(X) = (1,X,X^2,\dots,X^d)$ \alert{polynomiale Regression} von Grad $d$.

\uncover<2->{
  Auch die polynomiale Regression ist eine lineare Regression, denn\\
  \textcolor{darkgray}{Sie ist linear in den Gewichten, nicht in den Features}.
}
\end{frame}



\begin{frame}
  \frametitle{Ab jetzt ins Notebook}
  Im Notebook schauen wir uns \alert{Overfitten bei der polynomialen Regression} und dann ein \alert{Anwendungsbeispiel} mit dem Melbourne Housing Datensatz an, bei dem wir auch Scikit-Learn's \alert{Pipelines} kennen lernen.

  Im Notebook lernen wir noch ein weiteres zentrales Thema im maschinellen Lernen kennen, was wir bei der linearen Regression anwenden: \alert{Regularisierung}.

  \vspace{1cm}
  
\uncover<2->{
  Den Ansatz der Basisfunktionserweiterung kann man noch wilder treiben - nächstes Mal beschäftigen wir uns dann mit den \alert{Kernel-Methoden}.
 }
\end{frame}


\begin{frame}[separator]
  \frametitle{Vielen Dank für Ihre Aufmerksamkeit!}
\end{frame}

\end{document}

%%% Local Variables:
%%% mode: latex
%%% TeX-master: t
%%% End:
